% P16-C4: a complete finite sufficient front-door identification component.
\section{Finite front-door identification}
\label{sec:p16-front-door-identification}

A measured mediator can identify an intervention effect even when a
latent common cause prevents the measured variables from satisfying the
back-door criterion. Pearl's overview illustrates this construction with
\(W_3\) in Figure~4 and the composition in Eq.~(26)
\cite[sec.~3.3.2, p.~116]{Pearl2009CausalOverview}.
The theorem below supplies explicit sufficient graph conditions and a
complete finite derivation. It uses the structural intervention and
back-door results of
\cref{sec:p16-finite-causal-models,sec:p16-graphical-separation-adjustment}.

\subsection{Graph conditions and the mediated composition}

\begin{hbdefinition}{def:p16-front-door-data}{A single-treatment, single-mediator front-door criterion}
Let \(\mathcal M\) be a finite structural causal model with DAG
\(G=(N,E)\) and mutually independent disturbances. Fix three
distinct vertices: treatment \(t\), mediator \(h\), and response
\(r\). Write
\[
 X=V_t,\qquad M=V_h,\qquad Y=V_r.
\]
Here \(M\) denotes the mediator random variable; \(\mathcal M\)
denotes the whole structural model. Each of \(X,M,Y\) has its
specified finite nonempty value set \(S_t,S_h,S_r\).
All other endogenous vertices, including latent ones, remain in \(G\).

The mediator satisfies the \emph{front-door criterion} here when:
\begin{enumerate}
\item every directed path of positive length from \(t\) to \(r\)
      contains \(h\);
\item the empty set blocks every back-door path from \(t\) to \(h\);
\item \(\{t\}\) blocks every back-door path from \(h\) to \(r\).
\end{enumerate}
Back-door paths and blocking have their exact definitions in
\cref{def:p16-active-path,def:p16-back-door-data}. The first condition
is allowed to hold vacuously when there is no directed treatment-response
path. No statistical independence is inferred from an unblocked path.

For \(J\subseteq N\), let \(G_{\overline J}\) delete every
edge whose terminal vertex belongs to \(J\). Abbreviate
\(G_{\overline{\{t\}}}\) by \(G_{\overline t}\), and likewise
for \(h\). These are the graphs of equation-replacement interventions.
For scalar values \(x\in S_t,m\in S_h\), use the singleton
indexed-assignment convention of \cref{def:p16-back-door-data}, and put
\[
 M^x=V_h^{t=x},\qquad Y^x=V_r^{t=x},\qquad
 Y^m=V_r^{h=m},\qquad Y^{x,m}=V_r^{\{t,h\}=(x,m)}.
\]
In the last tuple, the coordinate at \(t\) is \(x\) and the
coordinate at \(h\) is \(m\). These superscripts label the
indicated interventions; they are not powers. The letters \(x,m\)
retain their treatment/mediator types even if the value sets overlap.
For named variables, \(Y^{X=x}\) and \(Y^{M=m}\) mean
\(Y^x\) and \(Y^m\), respectively.
\index{front-door criterion}\index{mediator!causal identification}
\end{hbdefinition}

\begin{hblemma}{thm:p16-front-door-graph-steps}{proved=annex:p16.front-door-graph-steps}{The needed properties of the modified graphs}
Under \cref{def:p16-front-door-data}:
\begin{enumerate}
\item \(t\notin\operatorname{de}_G(\{h\})\);
\item there is no directed path from \(t\) to \(r\) in
      \(G_{\overline h}\);
\item the empty set blocks every back-door path from \(h\) to
      \(r\) in \(G_{\overline t}\).
\end{enumerate}
The first conclusion makes \(\{t\}\) an eligible nondescendant
adjustment set for intervening at \(h\). The third conclusion is
a statement in the graph after fixing the treatment.
\end{hblemma}

\begin{hbtheorem}{thm:p16-front-door-composition}{proved=annex:p16.front-door-composition}{Composition through the mediator}
Under the same graph and model assumptions, for every \(x,m,y\)
of their specified types with \(\mathbb P(M^x=m)>0\),
\[
 \mathbb P(Y^x=y\mid M^x=m)=\mathbb P(Y^m=y).
\]
Consequently, for every \(x\in S_t,y\in S_r\),
\[
 \mathbb P(Y^x=y)
    =\sum_{m\in S_h}\mathbb P(M^x=m)\mathbb P(Y^m=y).
\]
The joint observational law need not be positive, and the original
event \(X=x\) may have zero probability. This identity concerns
the specified model's intervention laws. Its identification from
observation requires the next hypotheses.
\end{hbtheorem}

\subsection{The observable functional and its support conditions}

\begin{hbdefinition}{def:p16-front-door-overlap}{The treatment-specific overlap needed for the formula}
Fix \(x\in S_t\) with \(\mathbb P(X=x)>0\), and define the
observational support sets
\[
 \mathcal X_+=\{x'\in S_t:\mathbb P(X=x')>0\},\qquad
 \mathcal H_x=\{m\in S_h:\mathbb P(M=m\mid X=x)>0\}.
\]
Require
\[
 \mathbb P(M=m,X=x')>0
       \quad\text{for every }m\in\mathcal H_x,
                             \ x'\in\mathcal X_+.
\]
Thus every mediator value that contributes at the target treatment must
occur at every positive treatment stratum used to adjust its response.
This hypothesis is sufficient for the displayed conditional ratios;
no necessity claim for every possible identification method is made.
\end{hbdefinition}

\begin{hbtheorem}{thm:p16-finite-front-door-identification}{proved=annex:p16.finite-front-door-identification}{The finite front-door formula}
Suppose the model and graph satisfy \cref{def:p16-front-door-data},
and \(x\) satisfies \cref{def:p16-front-door-overlap}. Then, for
every \(y\in S_r\),
\[
\begin{aligned}
 \mathbb P_{\mathcal M}(Y=y\mid\operatorname{do}(X=x))
 ={}&\sum_{m\in\mathcal H_x}\mathbb P(M=m\mid X=x)\\[-2pt]
 &{}\cdot\sum_{x'\in\mathcal X_+}
       \mathbb P(Y=y\mid M=m,X=x')\mathbb P(X=x').
\end{aligned}
\]
All conditional probabilities in the expression have positive
denominators. Mediator values absent at the target treatment and
treatment strata of zero mass are omitted.

If \(\{t,h,r\}\subseteq O\) is observed, the right side is
determined by \(p_{\mathcal M,O}\). Hence it identifies the
intervention response law in the class with the stated graph, finite
value sets, independent disturbances and overlap. No observation of
the remaining endogenous vertices is required. Summing against any
specified function \(g:S_r\to\mathbb R\) gives its intervention mean; a contrast
of two treatment values requires the hypotheses at both values.
\end{hbtheorem}

The outer factor identifies the treatment-to-mediator law. The inner
sum identifies the mediator-to-response law by adjusting for treatment.
The composition theorem justifies combining them. Applying the two
adjustments alone would not justify that final composition when a
treatment-response path bypasses the mediator.

\subsection{A complete calculation with a latent common cause}

\begin{hbdefinition}{def:p16-front-door-binary-model}{A binary model with an unobserved root}
Name four vertices \(U,X,M,Y\), each with value set \(\{0,1\}\).
On \(\Omega=\{0,1\}^3\), assign outcome \((u,a,b)\) mass
\[
 q(u,a,b)=\tfrac12\alpha(a)\beta(b),\qquad
 \alpha(0)=\tfrac45,\quad\alpha(1)=\tfrac15,\quad
 \beta(0)=\tfrac34,\quad\beta(1)=\tfrac14.
\]
The disturbances are \(U_U=u\), \(U_X=a\), \(U_M=b\), and
the constant \(U_Y=0\), with total equations
\[
 U=U_U,\qquad X=U\oplus U_X,\qquad
 M=X\oplus U_M,\qquad
 Y=\mathbf1_{\{M=1\text{ or }U=1\}}.
\]
The allowed edges are \(U\to X\), \(X\to M\), \(M\to Y\),
and \(U\to Y\). The observed vertices are \(X,M,Y\), while
\(U\) is latent. Exclusive-or \(\oplus\) has its earlier bit
meaning. Independence here concerns the four disturbances; it does
not assert that the endogenous \(U\) and \(X\) are independent.
\end{hbdefinition}

\begin{hbtheorem}{thm:p16-front-door-binary-calculation}{proved=annex:p16.front-door-binary-calculation}{Observed association and the identified intervention law}
The model in \cref{def:p16-front-door-binary-model} satisfies all
three front-door graph conditions and overlap at both treatment values.
Its observational law on the measured coordinates is
\[
\begin{array}{cc|cc}
 X&M&Y=0&Y=1\\\hline
 0&0&12/40&3/40\\
 0&1&0&5/40\\
 1&0&1/40&4/40\\
 1&1&0&15/40
\end{array}
\]
Therefore
\[
 \mathbb P(Y=1\mid X=0)=2/5,\qquad
 \mathbb P(Y=1\mid X=1)=19/20.
\]
The front-door formula and direct equation replacement both give
\[
 \mathbb P(Y^{X=0}=1)=5/8,\qquad
 \mathbb P(Y^{X=1}=1)=7/8.
\]
Thus the observational response contrast is \(11/20\), whereas
the average intervention contrast is \(1/4\). The measured mediator
identifies this effect even though no subset of the other measured
coordinates satisfies the back-door criterion for \(X\) and \(Y\).
\end{hbtheorem}

\subsection{Why interception of every causal path matters}

\begin{hbdefinition}{def:p16-front-door-bypass-model}{A fully observed model with a direct bypass}
Let the vertices be \(X,M,Y\), all binary, with edges
\(X\to M\), \(M\to Y\), and \(X\to Y\). Let \(U_X\)
be a fair bit and \(U_M\) an independent bit with probability
\(1/4\) of being one; the disturbance for \(Y\) is constant.
Specify
\[
 X=U_X,\qquad M=X\oplus U_M,\qquad
 Y=\mathbf1_{\{X=1\text{ or }M=1\}}.
\]
All three vertices are observed. The two nonconstant disturbance
coordinates give an outcome space of pairs \((c,b)\in\{0,1\}^2\),
with \(U_X(c,b)=c\), \(U_M(c,b)=b\), and masses
\(\frac12\beta(b)\), using the \(\beta\) in
\cref{def:p16-front-door-binary-model}.
\end{hbdefinition}

\begin{hbtheorem}{thm:p16-front-door-bypass-counterexample}{proved=annex:p16.front-door-bypass-counterexample}{Removing the interception condition invalidates the formula}
The preceding model satisfies the second and third front-door graph
conditions and the stated overlap at both treatment values. It fails
the first condition because of \(X\to Y\). Evaluating the
observational front-door expression nevertheless gives \(5/8\) at
\(X=0\) and \(7/8\) at \(X=1\), while the actual intervention
response probabilities are \(1/4\) and \(1\), respectively.

The intervention response is still identified here by the empty-set
back-door formula \(\mathbb P(Y=y\mid X=x)\). Thus this
counterexample proves failure of the proposed front-door expression
when the criterion is removed; it does not prove universal
nonidentifiability of the effect.
\end{hbtheorem}

\begin{hbmethod}{def:p16-front-door-calculation}{Calculating a justified front-door functional}
Specify the complete graph and the observed treatment, mediator and
response. Check each of the three path conditions. For the desired
treatment value, form \(\mathcal X_+\) and \(\mathcal H_x\)
from the observational law and check the displayed joint-event
denominators. First compute the conditional mediator law, then the
response conditional on mediator and treatment. Average the latter
over the marginal treatment distribution, and average those results
over the target treatment's mediator distribution. The finite
identification theorem proves the resulting formula under exactly
these assumptions.
\end{hbmethod}

\Cref{sec:p16-do-calculus-soundness} next proves the three general
do-calculus transformation rules for the finite structural setting.
Completeness of graphical identification, multiple-mediator formulas
and general counterfactual identification require their additional
constructions. The numerical examples specify mathematical laws and
do not claim an empirical treatment effect or an unbiased finite-sample
estimator.
